% numerical_cross_validation.tex -- 市场数值证据
\section{数值结果、恒等式校核与验证边界}\label{sec:market-numerical}

本节保留既有 macOS Release 平台模型的历史验证结果；本次规则文档对齐未重跑 C++ 测试。
执行目标为 \srcpath{build/macos-release/tests/hacdcpf\_test\_market\_simulation}（目标名带
\texttt{hacdcpf\_} 前缀，以区别于 MIPSolvers full-dev 的同名 \texttt{test\_market\_simulation}
目标，见 \srcpath{tests/CMakeLists.txt:54-57}），记录为 22 cases、845 assertions。
这是市场模块内部回归证据，不是独立商业市场软件 oracle。
尤其不能作为南方区域 2025 年 V1.0 的 98 点 SCUC、分省备用、水库或独立 LMP 模型的数值证据。

\subsection{24 时段日前链路}
固定算例完成 24 个 SCUC 时段、24 个定价时段和 24 个 AC 认证时段，结果状态为
\texttt{converged}，\fld{num\_pricing\_converged}=24，\fld{num\_ac\_secure}=24。
算例含 3 台在役机组、9 条 AC 支路；估计 SCUC 变量数 2088、二元变量 72、SCED 变量数
1872，LODF 稠密存储估计 2456 bytes。某一代表时段的备用需求为 11.34 MW，备用出力和
需求差小于 $10^{-7}$ MW，负荷切除为 0 MW，所有 AC LMP 有限。

\begin{center}\small
\begin{tabular}{@{}L{6.0cm}rr@{}}
\toprule 指标 & 实测 & 门槛\\\midrule
定价收敛时段数 & 24/24 & 全部\\
AC 安全认证时段数 & 24/24 & 全部\\
最大备用平衡误差 & $<10^{-7}$ MW & $10^{-7}$ MW\\
负荷切除 & 0 MW & $<10^{-6}$ MW\\
SCUC warm start/结构提示/分支优先级 & true/true/true & 必须存在\\
\bottomrule
\end{tabular}
\end{center}

\subsection{交叉校核层级}
测试同时检查向量长度、备用平衡、LMP 有限性、结算记录和性能剖面；这些是源码结果
与契约的同实现一致性校核。第~\ref{sec:market-xref}~节的独立方程 oracle 现已在铜板定价
范围内提供独立重算的 LMP 与结算恒等式校核；含网阻塞 LMP、MATPOWER 市场 oracle 与跨引擎
全网 LMP 比较仍未覆盖，故除铜板切片外不能把 24/24 写成市场经济最优性证明。

\subsection{独立方程 oracle：铜板 SCED/LMP 与结算恒等式}\label{sec:market-xref}
注册测试 \srcpath{market\_sced\_cross\_validation} 以两个确定性单母线铜板算例把生产日前市场
\srcpath{src/market/market\_simulation.cpp:run\_day\_ahead\_market} 的定价与结算结果与一个不链接
hacdcpf 的独立 Python oracle 对照。证据生成器
\srcpath{tools/market\_validation/validate\_market\_xref.cpp:emit\_case} 用线性成本机组（边际
20 与 40 currency/MWh，各 100 MW）在负荷 60 与 140 MW 下清算；oracle
\srcpath{tools/market\_validation/run\_cross\_validation.py} 仅凭机组报价与负荷独立重算优先级调度栈
与边际（定价）机组，得到无损单母线上均一的 LMP，并逐项核对收益充分性恒等式：
资源电量收益 $=\sum_g \mathrm{LMP}_{bus(g)}p_g$、用户电量支付 $=\mathrm{LMP}\cdot D$、单母线阻塞租金为零、
现金流残差为零。

复现命令为
\begin{verbatim}
ctest --test-dir build/macos-release \
  -R '^market_sced_cross_validation$' --output-on-failure
\end{verbatim}

两个算例的解析期望为：负荷 60 MW 时仅廉价机组边际，$\mathrm{LMP}=20$、出力 $(60,0)$、用户
支付与资源收益均为 $1200$；负荷 140 MW 时廉价机组满发、调峰机组边际，$\mathrm{LMP}=40$、出力
$(100,40)$、用户支付与资源收益均为 $5600$；两例阻塞租金与现金流残差均为 $0$。准入阈值为 $10^{-6}$；
既有记录中九类检查最大误差为 $0$。负控制可复现：\srcpath{run\_cross\_validation.py} 的
\texttt{-{}-negative-control} 开关把第一个算例首条母线 LMP 人为改写 $+1.0$ 元/MWh 后再校验，
预期 oracle 检出并整体判 FAIL，此时脚本以退出码 0 报告“负控制如期失败”；若被改写的 LMP 反而
通过校验，脚本以退出码 1 报错（说明检查对 LMP 错误不敏感）。复现命令为
\begin{verbatim}
python3 tools/market_validation/run_cross_validation.py \
  --executable build/macos-release/tests/market_sced_xref \
  --out /tmp/market_xref_neg.json --markdown /tmp/market_xref_neg.md \
  --negative-control
\end{verbatim}
该 oracle 校核铜板经济
调度的边际定价与结算恒等式；有网络阻塞的 LMP、损耗、备用协优和 N-1 仍由市场 Catch2 套件验证，
本节不以铜板一致性冒充含网市场的外部准确度。

\subsection{失败与准入}
SCUC 达到相对 gap 但未证明最优时，结果必须保留 \fld{scuc\_optimality\_proven=false}；
AC 认证失败不重做商业出清，但 \fld{result.feasible} 置 \texttt{false} 且 \fld{status} 标记为
\fld{"ac\_validation\_failed"}（\srcpath{src/market/market\_simulation.cpp:5259-5269}），
同时 \texttt{SecurityResult} 标为失败；未支持资产通过
\fld{unsupported\_assets} 显式退回。任何报告若省略这三个字段，均不满足本手册准入。
